% =============================================================================
% Project : THE_PHILOSOPHY_OF_VIABILITY
% Book    : Book I — Ontology of Viability
% Book ID : VS.B01
%
% File    : 4_ASSETS/illustrations/
%           VS.B01.P01.C03.S02.SS08.P002.IL01.
%           differentiated_determinations_and_unity.tex
%
% Asset   : IL01 — Differentiated Determinations and Unity
% Section : 02 — Unity of Identity
% Chapter : 03 — Unity and Plurality
%
% Purpose :
%   To illustrate how differentiated elements may coexist
%   within one organized situation without implying that
%   their organization establishes one identity.
%
% Interpretive Limit :
%   The illustration does not prove the unity, singularity,
%   or number of identities involved in the depicted situation.
%
% Requires (preamble):
%   \usepackage{tikz}
%   \usetikzlibrary{calc, positioning, fit}
% =============================================================================

\begin{figure}[htbp]
    \centering

    \begin{tikzpicture}[
        font=\small,
        node distance=0.7cm,
        % --- base box ---
        box/.style={
            draw, rounded corners=3pt, align=center,
            inner sep=5pt, line width=0.7pt
        },
        % --- node roles ---
        detbox/.style={box, text width=2.6cm, minimum height=1.2cm,
                       fill=blue!8, draw=blue!55!black},
        orgbox/.style={box, text width=9.9cm, minimum height=1.0cm,
                       fill=gray!10, draw=black!60},
        % --- frame and labels (frame = scope of the example, not an
        %     identity boundary) ---
        vsgroup/.style={draw=black!45, dashed, rounded corners=6pt,
                        inner sep=10pt, line width=0.6pt},
        grouptitle/.style={font=\small\bfseries, align=center},
        legendtext/.style={font=\scriptsize},
        % --- undirected lines (no arrowheads: no order or priority) ---
        vsrel/.style={draw=black!70, line width=0.8pt},
        vslink/.style={draw=black!55, dashed, line width=0.7pt}
    ]

    % ---------------------------------------------------------------
    % Differentiated determinations
    % ---------------------------------------------------------------
    \node[detbox] (b) {\textbf{Determination B}};
    \node[detbox, left=of b]  (a) {\textbf{Determination A}};
    \node[detbox, right=of b] (c) {\textbf{Determination C}};

    \node[grouptitle, above=0.45cm of b] (ttl) {One Organized Situation};

    % ---------------------------------------------------------------
    % Relations and organization (shared field below)
    % ---------------------------------------------------------------
    \node[orgbox, below=1.0cm of b] (org) {
        \textbf{Relations and Organization}
    };

    \draw[vsrel] (a) -- (b);
    \draw[vsrel] (b) -- (c);

    \draw[vslink] (a.south) -- (a.south |- org.north);
    \draw[vslink] (b.south) -- (b.south |- org.north);
    \draw[vslink] (c.south) -- (c.south |- org.north);

    % ---------------------------------------------------------------
    % Frame: scope of the example
    % ---------------------------------------------------------------
    \node[vsgroup, fit=(ttl)(a)(b)(c)(org)] (frame) {};

    % ---------------------------------------------------------------
    % Note (outside the frame: identity is not part of the situation)
    % ---------------------------------------------------------------
    \node[below=0.5cm of frame, text width=10cm, align=center,
          font=\footnotesize\itshape] (note) {
        Differentiation and organization are both represented;
        identity remains an independent ontological question.
    };

    % ---------------------------------------------------------------
    % Legend
    % ---------------------------------------------------------------
    \node[below=0.35cm of note, legendtext] {
        \tikz[baseline=-0.5ex]\draw[vsrel] (0,0) -- (0.55,0);~relation among elements
        \quad\quad
        \tikz[baseline=-0.5ex]\draw[vslink] (0,0) -- (0.55,0);~element within the organizing field
    };

    \end{tikzpicture}

    \caption{Differentiated Determinations Within an Organized Situation.
    The elements are distinct and related within one depicted
    configuration. The frame represents the scope of the example,
    not a demonstrated boundary of one identity.}
    \label{fig:differentiated-determinations-and-unity}
\end{figure}

% =============================================================================
% END OF ASSET IL01 — SECTION 02 OF CHAPTER 03
% =============================================================================